\documentclass[../../General.tex]{subfiles}
\begin{document}

\section{Propuesta de solución}

Para superar las limitaciones de la instrumentación estática, esta propuesta plantea el diseño y desarrollo de un robot móvil especializado en la medición de variables climáticas en plantas fotovoltaicas. A diferencia de la infraestructura fija, este sistema funcionará como una plataforma móvil adaptada para transportar instrumentos meteorológicos, lo que le permitirá realizar un barrido continuo y focalizado en diversas secciones de la planta. Al desplazarse estratégicamente por el terreno, el vehículo recolectará datos de alta resolución que reflejen con mayor exactitud la dispersión geográfica de la irradiancia, el mosaico de microclimas y las condiciones ambientales locales. Esto tiene el fin de reducir la incertidumbre en la estimación de la generación de energía, en cumplimiento con los métodos de evaluación exigidos por las normativas internacionales. Para complementar su funcionalidad, el sistema almacenará y transmitirá la información capturada de forma inalámbrica, facilitando una plataforma de monitoreo remoto donde los operadores podrán visualizar el estado ambiental de la instalación en tiempo real.


\subsection{Objetivos generales}

\subsubsection{Objetivo correspondiente a Trabajo de Fin de Carrera 1 (1MTR01)}
Elaborar el diseño conceptual y el anteproyecto de un robot autónomo para la medición dinámica de variables climáticas en sistemas fotovoltaicos, orientado a caracterizar la formación de microclimas y variaciones climáticas en la topografía.

\subsubsection{Objetivo correspondiente a Trabajo de Fin de Carrera 2 (1MTR02)}
Desarrollar la ingeniería de detalle, elaborar el proyecto definitivo e implementar un prototipo virtual del robot autónomo para la medición de variables climáticas en sistemas fotovoltaicos, validando de forma integrada su comportamiento cinemático, control autónomo y arquitectura de telemetría.

\subsubsection{Objetivo para el proyecto general}
Desarrollar, dimensionar y validar integralmente un sistema mecatrónico autónomo móvil para la monitorización e inspección de variables climáticas en plantas fotovoltaicas emplazadas sobre terrenos irregulares, culminando en el proyecto definitivo y su validación mediante prototipado virtual.

\subsection{Objetivos específicos}

\subsubsection{Objetivos correspondientes a Trabajo de Fin de Carrera 1 (1MTR01)}
\begin{itemize}
    \item Determinar los requerimientos funcionales, cinemáticos y de diseño del sistema mecatrónico a partir del análisis de las severas condiciones topográficas y la dispersión climática en los parques solares.
    \item Desarrollar el diseño conceptual del sistema de locomoción, identificando sus funciones mecánicas y seleccionando los principios de tracción óptimos para garantizar un desplazamiento estable sobre terrenos irregulares.
    \item Desarrollar el diseño conceptual del mecanismo articulado, abstrayendo sus requerimientos cinemáticos y de actuación para asegurar la orientación independiente y precisa de los sensores climáticos.
    \item Estructurar la lógica del sistema de control autónomo, definiendo las funciones de procesamiento de señales, navegación y actuación requeridas para el seguimiento de rutas y la evasión de obstáculos.
    \item Plantear la arquitectura para el subsistema de telemetría y comunicaciones para garantizar la transmisión confiable de datos entre el robot y la central de control.
    \item Seleccionar el concepto de solución óptimo del sistema mediante una evaluación técnico-económica de las diferentes alternativas planteadas.
\end{itemize}

\subsubsection{Objetivos correspondientes a Trabajo de Fin de Carrera 2 (1MTR02)}
\begin{itemize}
    \item Realizar el diseño de detalle mecánico y estructural del sistema de tracción y de la base articulada para orientación de la instrumentación meteorológica, seleccionando materiales normalizados y generando la planimetría de fabricación correspondiente.
    \item Diseñar el sistema electrónico, instrumentación de variables ambientales y la arquitectura de alimentación energética independiente, dimensionando el almacenamiento en baterías y los circuitos de acondicionamiento de señales.
    \item Desarrollar e implementar los algoritmos de control de navegación autónoma, seguimiento de trayectorias predefinidas y evasión de obstáculos en el entorno operativo de la planta solar.
    \item Diseñar y configurar la infraestructura de comunicación inalámbrica y la interfaz gráfica de telemetría para la supervisión remota y recepción de datos climáticos en tiempo real.
    \item Construir y validar el prototipo virtual del sistema mecatrónico mediante simulaciones dinámicas y cinemáticas multicuerpo, pruebas en lazo cerrado de control y análisis de tráfico de red.
    \item Elaborar el proyecto definitivo del sistema mecatrónico, incluyendo la lista exhaustiva de materiales (BOM), el presupuesto detallado de implementación, el análisis de costos y la planificación de los protocolos de mantenimiento preventivo y correctivo.
\end{itemize}

\subsubsection{Objetivos para el proyecto general}
\begin{itemize}
    \item Diseñar a detalle una base de transporte con alta capacidad de tracción, a partir del concepto mecánico seleccionado, para operar de manera estable sobre los terrenos irregulares que caracterizan a los parques solares.
    \item Diseñar una base móvil articulada e integrada al vehículo, basándose en la arquitectura óptima definida previamente, capaz de moverse y orientarse de forma independiente para adaptarse a la posición óptima requerida por los diversos sensores ambientales.
    \item Desarrollar un algoritmo de control autónomo, implementando la lógica estructurada en el diseño conceptual, que proporcione al vehículo la inteligencia necesaria para seguir rutas de inspección predefinidas y reaccionar ante imprevistos mediante la evasión de obstáculos.
    \item Implementar una infraestructura de comunicación inalámbrica, materializando el esquema tecnológico propuesto, y definir la interfaz de telemetría y control para garantizar que los datos recolectados lleguen al centro de control sin interrupciones y permitiendo la intervención manual remota.
    \item Elaborar el proyecto definitivo del sistema incluyendo la elaboración de la lista de materiales, la estimación total de costos de implementación y la planificación de los protocolos de mantenimiento.
\end{itemize}

\end{document}

